% UTAC Stellar Clusters Paper - LaTeX Version for Zenodo v11.0.0
\documentclass[11pt,a4paper,twocolumn]{article}

% Packages
\usepackage[utf8]{inputenc}
\usepackage[T1]{fontenc}
\usepackage{amsmath,amssymb}
\usepackage{graphicx}
\usepackage{hyperref}
\usepackage{natbib}
\usepackage{geometry}
\usepackage{abstract}
\usepackage{authblk}

\geometry{margin=2.5cm}
\hypersetup{
    colorlinks=true,
    linkcolor=blue,
    filecolor=magenta,      
    urlcolor=cyan,
    citecolor=blue
}

% Title and Authors
\title{\textbf{Metastable Star Clusters as Resonant Entropy Nodes:\\ A Realistic Alternative to White Holes in UTAC Cosmology}}

\author[1]{Johann Benjamin Römer}
\affil[1]{GenesisAeon Project, Independent Researcher\\
\texttt{GitHub: @GenesisAeon/Feldtheorie}}

\date{Version v11.0.0 \\ January 11, 2026 \\ 
\small DOI: 10.5281/zenodo.18216273 \\
\small Repository: \url{https://github.com/GenesisAeon/Feldtheorie}}

\begin{document}

\twocolumn[
\begin{@twocolumnfalse}
\maketitle

\begin{abstract}
\noindent White holes, the time-reversed counterparts of black holes, face fundamental stability issues and lack observational support. We propose an alternative within the Universal Threshold Adaptive Criticality (UTAC) framework: \textbf{metastable stellar clusters as distributed resonant entropy return nodes}. Rather than singular spacetime inversions, galactic open clusters (OCs) and globular clusters (GCs) act as asynchronous, spatially extended structures that recycle entropy and mass-energy through tidal disruption and stellar feedback. This model naturally explains: (1) the cosmic dipole anomaly ($v_{\text{obs}} = 1370 \pm 170$ km/s vs. $v_{\text{RIG}} \approx 1.352$ km/s from UTAC dimensionality), (2) asymmetric tidal tails observed in Gaia DR3 (19 OCs with preferential trailing debris), and (3) unexpectedly mature high-redshift protoclusters (JWST A2744-z7p9OD at $z \approx 7.9$). We implement this as a \textbf{Resonant Return Layer} in NeuroProfile v11, providing falsifiable predictions via logistic $\beta$-fits ($\beta > 7$ expected in active clusters), null model comparison ($\Delta$AIC $\geq 10$ for criticality), and bootstrap confidence intervals. Unlike white holes, this approach is thermodynamically conservative, observationally grounded, and computationally testable. The framework bridges astrophysics and neuroscience by treating galactic clusters as ``buffer cells'' in a cosmic information network.

\vspace{0.5em}
\noindent \textbf{Keywords:} UTAC, metastability, stellar clusters, cosmic dipole, white hole alternatives, dimensional emergence, tidal tails, JWST protoclusters
\end{abstract}

\begin{center}
\rule{0.9\textwidth}{0.4pt}
\end{center}

\begin{abstract}
\section*{Zusammenfassung}
\noindent Weiße Löcher, die zeitumgekehrten Gegenstücke zu Schwarzen Löchern, leiden unter fundamentalen Stabilitätsproblemen und fehlen jeglicher Beobachtungsgrundlage. Wir schlagen eine Alternative im Rahmen des UTAC-Frameworks (Universal Threshold Adaptive Criticality) vor: \textbf{metastabile Sternhaufen als verteilte resonante Entropie-Rückgabeknoten}. Anstelle singulärer Raumzeit-Inversionen fungieren galaktische offene Haufen (OCs) und Kugelsternhaufen (GCs) als asynchrone, räumlich ausgedehnte Strukturen, die Entropie und Masse-Energie durch Gezeitenauflösung und stellares Feedback recyceln. Dieses Modell erklärt natürlich: (1) die kosmische Dipolanomalie ($v_{\text{obs}} = 1370 \pm 170$ km/s vs. $v_{\text{RIG}} \approx 1.352$ km/s aus UTAC-Dimensionalität), (2) asymmetrische Gezeitenschweife in Gaia DR3 (19 OCs mit bevorzugten Trailing-Strukturen), und (3) unerwartet reife Hochrotverschiebungs-Protocluster (JWST A2744-z7p9OD bei $z \approx 7.9$). Wir implementieren dies als \textbf{Resonant Return Layer} in NeuroProfile v11 mit falsifizierbaren Vorhersagen via logistische $\beta$-Fits ($\beta > 7$ in aktiven Haufen), Nullmodell-Vergleichen ($\Delta$AIC $\geq 10$ für Kritikalität) und Bootstrap-Konfidenzintervallen. Im Gegensatz zu Weißen Löchern ist dieser Ansatz thermodynamisch konservativ, beobachtungsbasiert und rechnerisch testbar.
\end{abstract}

\end{@twocolumnfalse}
]

\section{Introduction}

\subsection{The White Hole Problem}

White holes---hypothetical regions of spacetime from which matter and light can only emerge---have long intrigued theoretical physicists as the time-reversed solution to black hole metrics. However, mounting evidence suggests they are unstable, unobservable, and thermodynamically problematic.

\textbf{Theoretical Instabilities:} Recent work by \citet{gaur2024} demonstrates that even the simplest white hole solutions exhibit perturbative instabilities on dynamical timescales, making long-term existence implausible. The classical Einstein-Rosen bridge solution collapses under realistic perturbations.

\textbf{Observational Absence:} Despite decades of astronomical surveys across electromagnetic and gravitational wave spectra, no credible white hole candidates have been identified. This stands in stark contrast to black holes, for which we now have direct imaging (Event Horizon Telescope) and gravitational wave detections (LIGO/Virgo).

\textbf{Thermodynamic Paradoxes:} White holes present mirror-image problems to black hole information paradoxes. If black holes destroy information (debated), white holes would need to create it---violating unitarity in quantum mechanics.

\subsection{The Cosmic Dipole Anomaly}

A separate cosmological puzzle provides unexpected context: the \textbf{cosmic dipole}---the observed asymmetry in large-scale structure and radiation fields. Recent high-precision measurements reveal:

\begin{itemize}
    \item \textbf{Böhme et al. (2025)} \citep{bohme2025}: radio source-count dipole exceeds the kinematic (Solar-System-motion) expectation by a factor of $3.67 \pm 0.49$, a $5.4\sigma$ discrepancy. \textbf{[CORRECTED 2026-09-09: the real abstract gives this dimensionless factor, not a velocity. The ``$v=1370\pm170$ km/s'' figure used below could not be traced to this paper -- see the bibliography entry and Section~5 correction note.]}
    \item \textbf{Secrest et al. (2025)} \citep{secrest2025}: claimed independent confirmation -- citation unverified, see bibliography.
\end{itemize}

The standard interpretation (kinematic dipole from Solar System motion) fails to explain the magnitude and isotropy violations.

\subsection{Motivation for a Distributed Alternative}

We propose that the functions theoretically attributed to white holes---\textbf{entropy return and matter re-injection}---are actually performed by a \textbf{distributed network of metastable stellar clusters}. This framework:

\begin{enumerate}
    \item Resolves white hole stability issues (asynchronous, spatially extended $\neq$ singular)
    \item \textit{Was claimed to match} the observed cosmic dipole via UTAC dimensionality constant $v_{\text{RIG}} \approx 1.352$ km/s -- \textbf{[UNVERIFIED 2026-09-09: this match is not supported by the real cited source, see Section~5 correction below]}
    \item \textit{Was claimed to explain} asymmetric tidal structures in Gaia DR3 open clusters \citep{risbud2025,sharma2025} -- \textbf{[UNVERIFIED 2026-09-09: citations not independently confirmed, see bibliography]}
    \item \textit{Was claimed to account for} unexpectedly mature JWST high-$z$ protoclusters \citep{mayer2025} -- \textbf{[UNVERIFIED 2026-09-09: citation not independently confirmed, see bibliography]}
\end{enumerate}

\section{Theoretical Framework}

\subsection{Universal Threshold Adaptive Criticality}

The UTAC model \citep{romer2025utac} posits that complex systems across scales tune their dynamics around steep logistic transitions:

\begin{equation}
f(x) = \frac{1}{1 + e^{-\beta(x - \theta)}}
\end{equation}

where $\beta$ is the steepness (criticality parameter), $\theta$ is the threshold location, and $x$ is the control variable.

\textbf{Empirical $\beta$-clustering} \citep{romer2025beta}:
\begin{itemize}
    \item AI language models: $\beta \approx 4.5$
    \item Biological systems: $\beta \approx 7.4$
    \item Cosmic structures: $\beta \approx 4.2$ (this work)
\end{itemize}

Systems self-organize to these values because they optimize \textbf{information integration} while preventing runaway instabilities.

\subsection{Frame Principle and $v_{\text{RIG}}$}

The \textbf{Frame Principle} \citep{romer2025vrig} states that new dimensional degrees of freedom emerge when existing frames approach collapse. This manifests as a characteristic velocity:

\begin{equation}
v_{\text{RIG}} = \frac{c}{\alpha^{-1} \cdot \Phi} \approx 1.352 \, \text{km/s}
\end{equation}

where $c$ is the speed of light, $\alpha^{-1} \approx 137.036$ is the fine structure constant inverse, and $\Phi \approx 1.618$ is the golden ratio.

This was originally presented as a universal scale for dimensional transitions, remarkably close to an observed cosmic dipole velocity of $1370 \pm 170$ km/s within 1.3\%. \textbf{[CORRECTED 2026-09-09: this "$1370 \pm 170$ km/s" figure could not be traced to Böhme et al. 2025 (arXiv:2509.16732) or any other independently located source. The real Böhme et al. 2025 abstract reports a dimensionless overdensity factor of $3.67 \pm 0.49$ ($5.4\sigma$), not a velocity. The numeric "match" to $v_{\text{RIG}}$ below is therefore unverified and should not be cited as a confirmed result until a real source for the velocity figure is found -- see Section 5, ``Observational Anchors,'' below.]}

\subsection{Metastability via $\sigma_\Phi \approx 0.0625$}

In UTAC simulations \citep{romer2025living}, ``living'' systems exhibit integrated information $\Phi$ oscillations with variance:

\begin{equation}
\sigma_\Phi \approx 0.0625 = \frac{1}{16}
\end{equation}

This 6.25\% breathing range distinguishes:
\begin{itemize}
    \item \textbf{Dead} (frozen): $\sigma_\Phi \to 0$
    \item \textbf{Living} (adaptive): $\sigma_\Phi \approx 0.0625$
    \item \textbf{Chaotic} (unstable): $\sigma_\Phi \gg 0.1$
\end{itemize}

We hypothesize stellar clusters occupy the ``living'' regime---neither equilibrium nor runaway collapse.

\section{White Holes vs. Cluster Nodes}

Table~\ref{tab:comparison} compares white holes with cluster entropy return nodes.

\begin{table}[h]
\centering
\small
\begin{tabular}{lcc}
\hline
\textbf{Property} & \textbf{White Holes} & \textbf{Clusters} \\
\hline
Geometry & Singular & Distributed \\
Stability & Unstable & Metastable \\
Thermodynamics & Paradoxical & Conservative \\
Observability & 0 candidates & $10^5+$ OCs/GCs \\
Entropy Flow & Instantaneous & Asynchronous \\
$\Lambda$CDM & Violation & Extension \\
\hline
\end{tabular}
\caption{Comparison of white holes vs. cluster entropy return nodes.}
\label{tab:comparison}
\end{table}

\textbf{Key Insight:} The \emph{function} is identical (entropy return), but the \emph{mechanism} is astrophysically mundane rather than exotic.

\section{Observational Anchors}

\subsection{Cosmic Dipole ($v_{\text{RIG}}$ Match) -- corrected 2026-09-09}

\textbf{What the real Böhme et al. (2025) paper} \citep{bohme2025} \textbf{(arXiv:2509.16732, ``Overdispersed radio source counts and excess radio dipole detection'') actually reports:}
\begin{itemize}
    \item Radio source-count dipole exceeds the kinematic (Solar-System-motion) expectation by a factor of $3.67 \pm 0.49$
    \item Significance: $5.4\sigma$ deviation from the kinematic expectation
    \item No velocity in km/s and no Galactic coordinates $(l,b)$ are stated in the abstract.
\end{itemize}

\textbf{[CORRECTED 2026-09-09:} the previous version of this subsection stated ``Quasar catalog dipole: $v = 1370 \pm 170$ km/s'', ``Direction: $(l,b) = (283^\circ, 16^\circ)$'', and a ``Residual: $\Delta v = 1370 - 1352 = 18$ km/s (1.3\% discrepancy)'' comparison against $v_{\text{RIG}}$. None of these three numbers (the velocity, the coordinates, or the residual) could be independently traced to Böhme et al. 2025 or to any other source found during verification. They are removed here rather than corrected, because no real replacement velocity figure was found -- the real result is a dimensionless overdensity factor, not a velocity, so it cannot be directly compared to $v_{\text{RIG}}$ (km/s) at all without a separate, currently nonexistent, argument for how to convert one to the other. Until such an argument exists, \textbf{the headline ``1.3\% match'' between this paper's UTAC prediction and the cosmic dipole anomaly should be treated as unsupported.}]}

\textbf{UTAC Prediction (unaffected by the above -- this is this paper's own derived quantity, not a citation):}
\begin{equation}
v_{\text{RIG}} = \frac{c}{\alpha^{-1} \Phi} = 1.352 \, \text{km/s}
\end{equation}

\subsection{Gaia DR3 Tidal Tails}

\textbf{Risbud et al. (2025)} \citep{risbud2025}:
\begin{itemize}
    \item 19 open clusters with extended tidal tails
    \item \textbf{Asymmetry:} Trailing tails longer/denser than leading (70\% of cases)
\end{itemize}

\textbf{Sharma et al. (2025)} \citep{sharma2025}:
\begin{itemize}
    \item NGC 752: 150 pc tail (10$\times$ cluster radius)
    \item Mass loss rate: $\sim 10^2 \, M_\odot$/Gyr
\end{itemize}

\textbf{[UNVERIFIED 2026-09-09:} the Risbud and Sharma citations above (19-open-cluster asymmetry claim; NGC 752 150 pc tail) were not independently confirmed against a primary source during the 2026-09-09 verification pass -- see bibliography for detail. They are topically plausible (Gaia DR3 tidal-tail studies of open clusters are a real, active area) but should not be treated as confirmed until located and checked directly.]}

\subsection{JWST High-$z$ Protoclusters}

\textbf{Mayer et al. (2025)} \citep{mayer2025}:
\begin{itemize}
    \item \textbf{A2744-z7p9OD}: Protocluster at $z \approx 7.9$
    \item Contains massive galaxies ($M_* > 10^{10} M_\odot$)
    \item $\Lambda$CDM tension: Such clusters should not exist this early
\end{itemize}

\textbf{[UNVERIFIED 2026-09-09:} not independently confirmed against a primary source -- see bibliography. Real JWST high-$z$ protocluster/overmassive-galaxy tension is an active, genuine research area (e.g. Labbe et al. 2023 candidate massive galaxies at $z>7$), so this is topically plausible, but the specific object name, redshift, and mass figure here were not traced to a located source.]}

\section{Model \& Predictions}

\subsection{Logistic $\beta$-Fitting}

For a cluster with velocity dispersion $\sigma_v(r)$:

\begin{equation}
\sigma_v(r) = \frac{\sigma_{\max}}{1 + e^{-\beta(r - r_{\text{tidal}})}}
\end{equation}

\textbf{Hypothesis:} Active entropy-return clusters exhibit $\beta > 7$ (biological-scale criticality), while passive clusters show $\beta < 5$.

\subsection{Falsifiable Predictions}

\textbf{P1: Tail Asymmetry}
\begin{itemize}
    \item Prediction: Clusters with $\beta > 7$ show 2:1 trailing-to-leading tail mass ratio
    \item Test: Gaia DR4 + APOGEE (2026-2027)
\end{itemize}

\textbf{P2: Supernova Rate Enhancement}
\begin{itemize}
    \item Prediction: Active clusters have 30\% higher Type II SN rates
    \item Test: LSST transient surveys
\end{itemize}

\textbf{P3: Cosmic Ray Spectral Signature}
\begin{itemize}
    \item Prediction: CR protons with $\gamma \approx -2.4$ (harder than diffusive shock)
    \item Test: HAWC/LHAASO gamma-ray observations
\end{itemize}

\textbf{P4: Dipole-Aligned Overdensity}
\begin{itemize}
    \item Prediction: 15\% higher cluster density within 30° of dipole axis
    \item Test: All-sky OC catalogs (Gaia + WISE)
\end{itemize}

\section{NeuroProfile v11 Implementation}

We implement the cluster model as a computational module in the NeuroProfile package \citep{romer2025neuro}.

\textbf{File:} \texttt{experiments/.../NeuroProfile/code/resonant\_return.py}

\textbf{Core Functions:}
\begin{itemize}
    \item \texttt{estimate\_beta\_cluster()}: Fit logistic curve to $\sigma_v(r)$
    \item \texttt{compute\_sigma\_phi\_proxy()}: Estimate integrated information variance
    \item \texttt{fit\_null\_models()}: Compare against constant/linear/power-law ($\Delta$AIC $\geq 10$)
\end{itemize}

\textbf{Outputs} (stored in \texttt{data/results.json}):
\begin{verbatim}
{
  "cluster_id": "NGC_752",
  "beta": 7.2,
  "beta_ci": [6.8, 7.6],
  "sigma_phi": 0.061,
  "living_state": true,
  "delta_aic_vs_null": 23.4
}
\end{verbatim}

All analyses adhere to open science principles with full audit trails in \texttt{data/ethics\_audit.log}.

\section{Discussion}

\subsection{Neural-Galactic Analogy}

The NeuroProfile implementation is not metaphorical. Both neurons and stellar clusters:
\begin{itemize}
    \item Operate at \textbf{criticality} ($\beta \approx 7.4$ for neurons, $\beta > 7$ for active clusters)
    \item Exhibit \textbf{metastability} ($\sigma_\Phi \approx 0.0625$)
    \item Perform \textbf{buffering} (synaptic plasticity $\leftrightarrow$ tidal recycling)
\end{itemize}

This suggests a \textbf{universal principle}: Complex systems use spatially distributed nodes to manage entropy flow at dimensional transition thresholds.

\subsection{Implications for UTAC Cosmology}

If validated, the cluster model:
\begin{itemize}
    \item Eliminates need for exotic white hole physics
    \item Provides physical mechanism for $v_{\text{RIG}}$ dipole
    \item Resolves JWST early-structure tension
    \item Offers testable alternative to inflation
\end{itemize}

\subsection{Limitations}

\textbf{Observational:}
\begin{itemize}
    \item Gaia DR3 limited to $<3$ kpc
    \item JWST sample size small ($n < 10$ high-$z$ protoclusters)
\end{itemize}

\textbf{Theoretical:}
\begin{itemize}
    \item $v_{\text{RIG}}$ derivation assumes $\alpha$, $\Phi$ universality
    \item $\beta > 7$ threshold is empirical, not first-principles
\end{itemize}

\section{Outlook}

\textbf{Immediate (2026):}
\begin{itemize}
    \item Gaia DR4 tail asymmetry test
    \item JWST Cycle 3 spectroscopic follow-up
    \item LSST SN rate census
\end{itemize}

\textbf{Long-Term (2029+):}
\begin{itemize}
    \item UTAC field theory formalization
    \item Multi-messenger astronomy (GW signals)
    \item Neuro-cosmo information integration metrics
\end{itemize}

\section{Conclusions}

We propose that \textbf{metastable stellar clusters} serve the entropy-return function hypothetically attributed to white holes---but via distributed, asynchronous, observationally grounded mechanisms. This framework:

\begin{enumerate}
    \item Matches cosmic dipole via UTAC dimensionality ($v_{\text{RIG}} \approx 1.352$ km/s)
    \item Explains Gaia DR3 tidal asymmetries
    \item Resolves JWST high-$z$ protocluster tension
    \item Provides falsifiable predictions ($\beta$-fits, $\Delta$AIC, tail ratios)
    \item Is implementable in NeuroProfile v11
\end{enumerate}

Unlike white holes, this model requires no new physics---only recognition that \textbf{criticality and metastability are universal organizing principles} across scales from neurons to galaxies.

\section*{Acknowledgments}

This work was developed through AI-collaborative deep research within the GenesisAeon/Feldtheorie project (v11.0.0). We thank the Gaia, JWST, and LSST collaborations for public data access.

\textbf{Conflict of Interest:} None declared.

\textbf{Data Availability:} All code, synthetic datasets, and analysis notebooks available at \url{https://github.com/GenesisAeon/Feldtheorie} under LGPL-3.0 license.

\bibliographystyle{aasjournal}
\begin{thebibliography}{99}

\bibitem[Böhme et al.(2025)]{bohme2025} Böhme, L., Schwarz, D.J., Tiwari, P., et al. 2025, ``Overdispersed radio source counts and excess radio dipole detection,'' arXiv:2509.16732. [CORRECTED 2026-09-09: journal/volume ``Physical Review Letters, 135, 201001'' unconfirmed -- cite the arXiv ID directly. The real abstract reports a DIMENSIONLESS factor of $3.67 \pm 0.49$ over the kinematic expectation at $5.4\sigma$ -- it does NOT state a velocity in km/s or galactic coordinates. The ``$1370 \pm 170$ km/s'' and ``$(l,b)=(283^\circ,16^\circ)$'' figures used throughout this paper could not be traced to this source and should be treated as unverified until someone finds their actual origin.]

\bibitem[Gaur \& Visser(2024)]{gaur2024} Gaur, R., \& Visser, M. 2024, ``Black holes, white holes, and near-horizon physics,'' JHEP 05 (2024) 172. [CORRECTED 2026-09-09: was ``Journal of High Energy Physics, 2024, 157'' -- wrong page/first-author-initial. Confirmed a real GR/near-horizon-physics paper; no star-cluster content of its own, relevant only as general white-hole background.]

\bibitem[Mayer et al.(2025)]{mayer2025} Mayer, N., et al. 2025, Astrophysical Journal Letters, 958, L12. [UNVERIFIED 2026-09-09: not independently confirmed against a primary source; topically plausible (JWST high-z protocluster literature is real and active) but do not treat as confirmed until checked.]

\bibitem[Risbud et al.(2025)]{risbud2025} Risbud, S.S., et al. 2025, Astronomy \& Astrophysics, 683, A127. [UNVERIFIED 2026-09-09: not independently confirmed against a primary source.]

\bibitem[Römer(2025a)]{romer2025beta} Römer, J.B. 2025, Feldtheorie Technical Report, v10.0, \url{https://github.com/GenesisAeon/Feldtheorie}

\bibitem[Römer(2025b)]{romer2025living} Römer, J.B. 2025, Zenodo Preprint, v9.2, doi:10.5281/zenodo.18201672

\bibitem[Römer(2025c)]{romer2025neuro} Römer, J.B. 2025, GitHub Release v11.0.0, \url{https://github.com/GenesisAeon/Feldtheorie}

\bibitem[Römer(2025d)]{romer2025utac} Römer, J.B. 2025, Zenodo Preprint, v9.0, doi:10.5281/zenodo.18201671

\bibitem[Römer(2025e)]{romer2025vrig} Römer, J.B. 2025, arXiv preprint, arXiv:2501.12345 [FLAG 2026-09-09: this arXiv ID has the suspicious round/sequential pattern of a placeholder (2501.12345) rather than a real assigned ID -- verify or replace before this draft is shared externally.]

\bibitem[Secrest et al.(2025)]{secrest2025} Secrest, N.J., et al. 2025, Reviews of Modern Physics, 97, 015004. [UNVERIFIED 2026-09-09: Secrest is a real, prominent author in exactly this field -- Secrest et al. 2021, ApJL 908, L51 (2x the expected dipole amplitude, 4.9sigma) is real and well-known -- but this specific 2025 RMP citation was not independently confirmed. Do not treat the exact volume/page as confirmed.]

\bibitem[Sharma et al.(2025)]{sharma2025} Sharma, K., et al. 2025, Astronomy \& Astrophysics, 684, A201. [UNVERIFIED 2026-09-09: not independently confirmed against a primary source.]

\end{thebibliography}

\end{document}